\section{Metodologia de simulacion M/M/1}

\subsection{Diseno experimental}

La unidad de tiempo elegida para el modelo de colas es la hora. Se fija una tasa de servicio base:

\begin{equation}
    \mu = 10 \text{ clientes/hora}.
\end{equation}

Las tasas de arribo se definen como fracciones de $\mu$:

\begin{equation}
    \lambda \in \{0.25\mu, 0.50\mu, 0.75\mu, 1.00\mu, 1.25\mu\}.
\end{equation}

Por lo tanto, los valores de $\rho$ evaluados son $0.25$, $0.50$, $0.75$, $1.00$ y $1.25$. Para cada valor se simula una cola infinita y colas finitas con:

\begin{equation}
    B \in \{0,2,5,10,50\}.
\end{equation}

El horizonte de cada corrida es de $10000$ horas y se ejecutan $10$ replicas por experimento. Las semillas se derivan de una semilla base registrada, para que los resultados sean reproducibles.

\subsection{Logica de eventos discretos}

La simulacion mantiene un reloj y procesa dos tipos de eventos:

\begin{itemize}
    \item llegada de un cliente;
    \item finalizacion de un servicio.
\end{itemize}

En cada llegada se agenda el siguiente arribo. Si el servidor esta libre, el cliente inicia servicio inmediatamente. Si el servidor esta ocupado y existe lugar en la cola, el cliente espera. Si la cola esta llena, el cliente se rechaza y se incrementa el contador de denegaciones.

Los tiempos entre arribos y los tiempos de servicio se generan con variables exponenciales:

\begin{align}
    T_a &\sim \mathrm{Exp}(\lambda), \\
    T_s &\sim \mathrm{Exp}(\mu).
\end{align}

\subsection{Metricas registradas}

Las medidas temporales se calculan integrando el estado del sistema a lo largo del horizonte simulado. Si $T$ es el tiempo total de simulacion:

\begin{align}
    \widehat{L} &= \frac{1}{T}\int_0^T N(t)\,dt, \\
    \widehat{L_q} &= \frac{1}{T}\int_0^T Q(t)\,dt, \\
    \widehat{U} &= \frac{\text{tiempo ocupado del servidor}}{T}.
\end{align}

Para cada cliente atendido tambien se registra el tiempo de espera en cola y el tiempo total dentro del sistema. A partir de esos datos se estiman:

\begin{align}
    \widehat{W} &= \text{tiempo promedio en sistema}, \\
    \widehat{W_q} &= \text{tiempo promedio en cola}.
\end{align}

La probabilidad de denegacion se estima como:

\begin{equation}
    \widehat{P}_{den} =
    \frac{\text{clientes rechazados}}{\text{clientes generados}}.
\end{equation}

\subsection{Resumen estadistico}

Para cada combinacion de parametros se calcula la media de las $10$ replicas, el desvio estandar muestral y un intervalo de confianza del $95\%$:

\begin{equation}
    \bar{X} \pm t_{0.975,9}\frac{s}{\sqrt{10}}.
\end{equation}

Cuando existe valor teorico estacionario, se reporta ademas el error porcentual:

\begin{equation}
    \text{Error \%} =
    100 \frac{\bar{X} - X_{\mathrm{teorico}}}{X_{\mathrm{teorico}}}.
\end{equation}

Los experimentos de cola infinita con $\rho \geq 1$ se conservan porque forman parte de las tasas pedidas, pero se etiquetan como casos sin regimen estacionario. En consecuencia, sus metricas se interpretan como resultados de horizonte finito y no como aproximaciones a una teoria estacionaria.
